
\problem{1}
Существует ли такое натуральное $n$, что число $3^n+2\cdot 17^n$ является точным квадратом?

\problem{2}
Натуральные числа $a$, $b$ и $c$ таковы, что числа $ab + 9b + 81$ и $bc + 9c + 81$ делятся на 2023. Докажите, что $ca + 9a + 81$ так же делится на 2023.

\problem{3}
При каких нечетных $s$ существуют такие целые $m$ и $n$, что $m^6 - 4n^3 + 3^s = 0$?

\problem{4}
При каком наименьшем натуральном $n$ существуют такие целые числа $a_1$, $a_2$, …, $a_n$, что уравнение $$x^2-2(a_1+a_2+\ldots+a_n)^2x + (a_1^4+a_2^4+\ldots+a_n^4+1)=0$$ имеет по крайней мере один целый корень?

\problem{5}
В ряд написан 200 натуральных чисел. Известно, что каждое число начиная со второго, либо в 9 раз больше предыдущего, либо в 2 раза меньше предыдущего. Может ли сумма всех этих чисел равняться $2022^{2022}$?

\problem{6}
Докажите, что если целые числа $x$ и $y$ таковы, что $x^2+xy+y^2$ делится на простое число $p$ вида $3k+2$, то и сами числа $x$ и $y$ делятся на $p$.

\problem{7}
Дано нечетное простое число $p$. Чему может быть равен НОД чисел $a+b$ и $\cfrac{a^p+b^p}{a+b}$, если $a$ и $b$ — взаимно простые целые числа?

\problem{8}
Пусть $p$ — нечетное простое число. Докажите, что $\sum\limits_{k=1}^{p-1} k^{2p-1}\equiv \frac{p(p+1)}{2} \ (mod \ p^2)$.

\problem{9}
Докажите, что для любого простого числа $p$ найдется такое натуральное число $n$, что число $6^n+3^n+2^n-1$ делится на $p$.

\problem{10}
Докажите, что для любого простого $p>3$ числитель несократимой дроби, равной $$1+\frac{1}{2^2}+\frac{1}{3^2}+\frac{1}{4^2}+\ldots+\frac{1}{(p-1)^2},$$
делится на $p$.

\problem{11}
Найти все приведенные целочисленные многочлены $Q(x)$, обладающих следующим свойством: существует натуральное $N$, такое, что для любого простого числа $p>N$ число $2\cdot Q(p)!+1$ делится на $p$.